% ===== 原文 pp.39–40 =====
% 译者 tr-06：§5 开头、§5.1 群同胚（含 5.1.1）；编号区间：式 (5.1)、式 (5.2)、命题 5.3
% 依据：split/p039.txt、split/p040.txt（仅用于定位）＋ pages/p039.png、p040.png（200 dpi，数学逐符号核对）

% [原文 p.39]
\section{Higgs 丛与对应关系}

\begin{quote}
\kaishu
近来，许多长期悬而未决问题最引人注目的解决，都来自数学不同领域思想出人意料的组合。

\rmfamily\hfill Sir Michael Atiyah
\end{quote}

在考虑复群的 Higgs 丛模空间时，人们可以追问：它与其他模空间（Higgs 丛的模空间，或其他数学对象的模空间）之间有哪些对应关系？它的哪些子空间具有有趣的对应与解释？下文我们将描述其中若干项，既涉及 Higgs 丛的整个模空间，也涉及第 4 章所定义的 Higgs 丛的膜。

% [Tutor Note: p.39 原文 "In what fallows we shall describe a few of these..." 中 "fallows" 应为 "follows"（原文笔误；同一措辞的正确拼写见 p.40 "In what follows, we shall give an overview..."）。]

\subsection{群同胚}

% [Tutor Note: p.39 原文小节标题作 "5.1. Group homeomorphisms."，但本节（含 5.1.1）通篇讨论 Lie 群之间的同态（homomorphism）与同源（isogeny），式 (5.1) 的 $\Psi : G_{\C}\to G'_{\C}$ 正是群同态。故 "homeomorphisms"（同胚）应为 "homomorphisms"（同态）；译文照原样作「群同胚」，正确说法为「群同态」。]

给定固定的黎曼面 $\Sigma$ 以及两个 Lie 群之间的同态
\begin{equation}
\Psi : G_{\C}\to G'_{\C}, \tag{5.1}
\end{equation}
则自然地诱导出表示空间之间的映射
\[
\Psi : \Rep(\Sigma,G_{\C})\to \Rep(\Sigma,G'_{\C}),
\]
其中 $\Rep(\Sigma,G)$ 是表示模去共轭作用所得的空间。正如讲义开头所提到的，曲面群的表示与 Higgs 丛之间存在对应，因此 $G_{\C}$-Higgs 丛的模空间之间也应有类似诱导的态射。于是可以追问：由群同态能得到 Higgs 丛之间的哪些对应？这些映射中又有哪些能下降为子群之间的对应？

% [Tutor Note: p.39 "the space of representations modulo conjugation" 指表示按共轭分类（即 $\Hom(\pi_1(\Sigma),G)$ 模去 $G$ 的共轭作用），并非把矩阵的共轭转置等同起来。可核对：式 (5.2) 的像 $P_{G_{\C}}\times_{\Psi}G'_{\C}$ 是 associated bundle，说明对应是在 $G_{\C}$-丛的层面发生的。]

% [Tutor Note: p.39 "which of those maps descend to correspondences between subgroups" 中的 "subgroups" 疑为 "subspaces" 之误，存疑待考：本段开头问的正是模空间的子空间（"whether any of its subspaces have interesting correspondences and interpretations"），§5 余下讨论的也都是模空间的子空间（膜）之间的对应。译文按原词照译为「子群」，未作改动。]

通过定义 2.9 可以看出，式 (5.1) 中的映射诱导出主 $G_{\C}$-Higgs 丛 $(P_{G_{\C}},\Phi)$ 上的映射
\begin{equation}
\Psi : (P_{G_{\C}},\Phi)\mapsto(P_{G_{\C}}\times_{\Psi}G'_{\C},\ \dd\Psi(\Phi)), \tag{5.2}
\end{equation}
其中 $P_{G'_{\C}} := P_{G_{\C}}\times_{\Psi}G'_{\C}$，而 Higgs 场的像是 $\Phi' := \dd\Psi\Phi$，这里
\[
\dd\Psi : \mathrm{ad}(P_{G_{\C}})\to\mathrm{ad}(P'_{G_{\C}})
\]
由映射 $\Psi$ 在单位元处的微分给出。

\para{5.1.1.\ 由同源诱导的映射}

回应上述问题的第一步，是考虑群同态由 Lie 群的同源（isogeny）给出的情形，下面即将看到这一点。通过低秩 Lie 代数之间的对应，可以发现不同 Higgs 丛之间的一些有趣关系。复 Lie 群之间的部分低秩同源为
\begin{align*}
\SO(3,\C) &\cong\Sp(2,\C)\cong\SL(2,\C);\\
\SO(4,\C) &\cong\SL(2,\C)\times\SL(2,\C);\\
\SO(5,\C) &\cong\Sp(4,\C);\\
\SO(6,\C) &\cong\SL(4,\C).
\end{align*}
% [Tutor Note: p.39 原文 "Some of the low rank isogenies between complex are" 在 "complex" 之后漏掉名词，应为 "between complex Lie groups are"：p.40 首句与之对举的是 "maps between real Lie algebras"，且所列四条同源都是复 Lie 群之间的。译文按其确切语义补出「Lie 群」。]
% [Tutor Note: p.39 四条低秩同源可用 Dynkin 图核对：$\SO(4,\C)\cong\SL(2,\C)\times\SL(2,\C)$ 即 $D_2=A_1\times A_1$；$\SO(6,\C)\cong\SL(4,\C)$ 即 $D_3=A_3$；$\SO(5,\C)\cong\Sp(4,\C)$ 即 $B_2=C_2$。首行 $\Sp(2,\C)$ 取 $2\times2$ 辛群，于是 $\Sp(2,\C)=\SL(2,\C)$，与 $\SO(3,\C)$ 同为复维数 3。]
% [原文 p.40]
这些对应关系诱导出实 Lie 代数之间的映射，可描述如下：
\begin{alignat*}{2}
\text{(a)}\ \fso(2,1) &\cong\fsl(2,\R), &\qquad \text{(f)}\ \fso(3,3) &\cong\fsl(4,\R),\\
\text{(b)}\ \fso(2,2) &\cong\fsl(2,\R)\times\fsl(2,\R), &\qquad \text{(g)}\ \fso(2,4) &\cong\fsu(2,2),\\
\text{(c)}\ \fso^*(4) &\cong\fsu(2)\times\fsl(2,\R), &\qquad \text{(h)}\ \fso(1,5) &\cong\fsu(4),\\
\text{(d)}\ \fso(1,3) &\cong\fsl(2,\C), &\qquad \text{(i)}\ \fso(2,3) &\cong\fsp(4,\R).\\
\text{(e)}\ \fso(1,4) &\cong\fsp(2,2),
\end{alignat*}

% [Tutor Note: p.40 第 (h) 条原文作 $\fso(1,5)\cong\fsu(4)$，应为 $\fso(1,5)\cong\fsu^*(4)$：$\fsu(4)$ 是紧 Lie 代数，而 $\fso(1,5)$ 非紧（二者复维数同为 15、秩同为 3，但紧性不同，故不可能同构）。旁证：同表第 (c) 条 $\fso^*(4)\cong\fsu(2)\times\fsl(2,\R)$ 就正确地带了星号。译文照原样排 $\fsu(4)$。]

(b)、(f) 中秩 2 与秩 3 的同源诱导出膜之间的映射，已在 [29] 中借助 Hitchin 纤维化的正则纤维作了描述；而 (c)、(d) 中奇异纤维内部膜之间的对应则由 [97] 处理。其余所有对应都尚未通过谱数据研究过，因此它们构成一份内涵丰富的、可据以提问的开放问题清单。下文我们将依据 [29] 与 [97] 概述这一领域以及人们可以提出的问题。关于进一步的背景，读者可参阅 [7]，其中对用 Higgs 丛的语言解释同源给出了更多说明。

由式 (5.2) 不难看出，向量丛与 Higgs 场层面的 Higgs 丛之间的对应可以自然地导出（见 [7,29]）。回忆一下：一个 $\SL(n,\C)$-Higgs 丛可以看作一对 $(E,\Phi)$，其中 $E$ 是 $\Sigma$ 上秩为 $n$、行列式平凡的全纯丛，而 $\Phi$ 是无迹的 Higgs 场。此外，一个 $\SO(n,\C)$-Higgs 丛可以写成一对 $(E,\Phi)$，其中 $E$ 是 $\Sigma$ 上秩为 $n$ 的全纯丛，带有正交结构 $Q$ 以及其行列式丛的一个相容平凡化——这里平凡化 $\delta : \det(E)\simeq\cO_{\Sigma}$ 与 $Q$ 相容，是指 $\delta^2$ 与由判别式 $Q : \Lambda^n E\to\Lambda^n E^*$ 给出的 $(\Lambda^n E)^2$ 的平凡化一致；并且 Higgs 场 $\Phi$ 满足
\[
Q(u,\Phi v) = -Q(\Phi u, w).
\]

% [Tutor Note: p.40 末式 "$Q(u,\Phi v) = -Q(\Phi u, w)$" 中的 $w$ 应为 $v$：这是 $\Phi\in\fso(n,\C)$ 关于对称双线性型 $Q$ 的反自伴性（即 $\Phi$ 落在 $\fso(n,\C)$ 中）的定义式，两端须用同一对向量 $(u,v)$；写成 $w$ 会使右端出现前文未定义的第四个向量。译文照原样排 $w$。]
% [Tutor Note: p.40 "an $\SL(n,\C)$-Higgs bundle can be seen as is a pair" 多出一个 "is"（应为 "can be seen as a pair"）；译文按正常语义处理，不照搬这一赘词。]

\begin{prop}{5.3}
由
\begin{align*}
\cI_2 :\ &\SL(2,\C)\times\SL(2,\C)\to\SO(4,\C);\\
\cI_3 :\ &\SL(4,\C)\to\SO(6,\C);
\end{align*}
给出的秩 2 与秩 3 的 Lie 群同源，诱导出 $\SL(2,\C)$-Higgs 丛与 $\SL(4,\C)$-Higgs 丛上的映射，它们分别由
\begin{align*}
\cI_2 :\ &(E_1,\Phi_1),(E_2,\Phi_2)\mapsto[E_1\otimes E_2,\ \Phi_1\otimes 1+1\otimes\Phi_2]\in\cM_{\SO(4,\C)},\\
\cI_3 :\ &(E,\Phi)\mapsto\bigl[\Lambda^2 E,\ \Phi\otimes 1+1\otimes\Phi\bigr]\in\cM_{\SO(6,\C)}.
\end{align*}
给出。
\end{prop}

\para{证明}把式 (5.1) 与 (5.2) 应用于相应的 $G_{\C}$-Higgs 丛，命题直接可得。\hfill$\square$

我们感兴趣的是通过谱数据来理解命题 5.3 中所描述的映射。为此，需要研究 Higgs 场特征值之间的对应——它给出谱曲线之间的对应；以及特征空间之间的对应——它给出谱曲线上线丛之间的对应。
